\section*{Abstrakt}\label{Abstrakt}

% Wymagane elementy:
% + Cel pracy
% + Postawione hipotezy
% + Materiał i metody badawcze
% +/- Najważniejsze wyniki
% +/- Podsumowanie i wnioski

% Wymagania:
% - Po polsku i po angielsku
% - We formie bezosobowej lub równoważników zdań

Celem pracy jest ocena znajomości i częstości stosowania suplementów diety wśród studentów, ze szczególnym uwzględnieniem najczęściej wybieranych preparatów, powodów ich stosowania oraz źródeł wiedzy.

Określenie częstości stosowania suplementów diety wśród studentów.
Identyfikacja najczęściej stosowanych suplementów.
Analiza powodów suplementacji.
Określenie głównych źródeł informacji.
Ocena poziomu wiedzy studentów na temat suplementów.
Analiza zależności między poziomem wiedzy a praktyką suplementacyjną.

Jako hipotezy przyjęto:
1. (Częstość) Studenci często stosują suplementy diety.
2. (Najpopularniejsze suplementy) Studenci najczęściej stosują suplementy kompozytowe oraz witaminę C.
3. (Powody) Studenci suplementują ze względu na stres związany ze studiami celem zadbania o zdrowie.
4. (Źródła) Głównym źródłem informacji jest internet.
5. (Wiedza) Studenci najczęściej mają niewystarczającą wiedzę na temat suplementów.
6. (Zależność) Studenci często stosują suplementy diety w sposób nieprawidłowy lub niebezpieczny.

Dane wykorzystane w badaniu zebrano za pośrednictwem ankiety internetowej udostępnionej wśród grup studenckich online.

W wyniku analizy zebranych danych można było stwierdzić zależność pomiędzy poziomem wiedzy a właściwą praktyką suplementacyjną.

Im wyższy poziom rzeczywistej wiedzy tym bardziej odpowiednie praktyki suplementacyjne.

% Słowa kluczowe

Dieta, Praktyka, Studenci, Suplementacja, Wiedza

